\bookchapter{第五十二章：四大元件與 Activity 生命週期}

\chapterintro{如果說 Compose 是 App 的“臉”，那這一章講的就是 App 的“骨架”。安卓系統有四大基本元件，所有 App 都是由它們拼成的；而其中最重要的 Activity，有一張每個安卓開發者都必須刻在腦子裡的“生命週期圖”。本章學完，你就能理解“為什麼切個後臺回來，我的 App 狀態就沒了”這類千古之謎。}

\booksection{一、四大元件：安卓的“四梁八柱”}

安卓官方定義了四種基本元件，分別負責一類工作：

\begin{center}
\begin{tabularx}{\linewidth}{|X|X|X|}
\hline
\textbf{元件} & \textbf{職責} & \textbf{打個比方} \tabularnewline
\hline
Activity & 一個可見的介面 & 店鋪的門面——使用者能看見、能互動 \tabularnewline
Service & 後臺幹活的"隱形人" & 倉庫工人——放音樂、下載檔案，不需要介面 \tabularnewline
BroadcastReceiver & 接收系統廣播 & 傳達室——"電池沒電了""網路斷了"的喇叭 \tabularnewline
ContentProvider & 跨 App 共享資料 & 公共資料櫃——比如把聯絡人共享給別的 App \tabularnewline
\hline
\end{tabularx}
\end{center}

\begin{center}\small\color{gray}安卓四大元件速覽\end{center}

新手階段你 95\% 的時間都在跟 Activity 打交道；Service 和 BroadcastReceiver 做音樂播放、訊息推送時會用到；ContentProvider 用得最少（系統已把常用資料封好了）。另外注意：四大元件都要在 AndroidManifest.xml 裡“登記戶口”，不登記系統就找不到它們。

\booksection{二、Activity 生命週期：安卓開發的第一張“藏寶圖”}

手機上的 App 隨時可能被打斷：來電話了、使用者切去回微信、系統記憶體不夠把你殺了……Activity 的一生，就是在這張圖上迴圈旅行：

onCreate（出生：初始化）

onStart（可見了）

onResume（前臺互動中）

onPause（部分被遮住）

onStop（完全看不見）

onDestroy（銷燬）

onRestart / onStart / onResume：回到前臺時原路返回

\begin{center}\small\color{gray}圖 4：Activity 生命週期（前進方向：上排出生 下排離場，回來時原路折返）\end{center}

\begin{center}
\begin{tabularx}{\linewidth}{|X|X|X|}
\hline
\textbf{回撥} & \textbf{什麼時候觸發} & \textbf{典型用途} \tabularnewline
\hline
onCreate & Activity 建立時（一次） & 載入介面、初始化資料 \tabularnewline
onStart & 介面變為可見 & 開始一些輕量邏輯 \tabularnewline
onResume & 拿到焦點、可互動 & 開始動畫、還原播放 \tabularnewline
onPause & 被部分遮擋/切走瞬間 & 暫停動畫、儲存草稿 \tabularnewline
onStop & 完全不可見 & 釋放重資源、暫停下載 \tabularnewline
onDestroy & 被銷燬前 & 最後的清理 \tabularnewline
\hline
\end{tabularx}
\end{center}

\begin{center}\small\color{gray}六個生命週期回撥的典型用途\end{center}

\textbf{避坑提醒：螢幕旋轉會“殺死”Activity}\par

新手最經典的翻車現場：App 轉一下螢幕，輸入框裡的字全沒了。原因：旋轉螢幕預設會銷燬並重建 Activity（走一遍 onDestroy onCreate）。解法：用 rememberSaveable 儲存狀態（Compose 裡替代 remember），它會在重建時自動把值找回來。第 6、7 章會反覆用到。

\booksection{三、Intent：App 裡外通行的“車票”}

Intent（意圖）是元件之間通訊的“車票”。你想從一個頁面跳到另一個頁面，或者開啟瀏覽器、調起相機，都要“買票上車”：

\begin{lstlisting}[language=Kotlin]
// 方式一：顯式 Intent——明確指定去哪個頁面
val intent = Intent(this, DetailActivity::class.java)
intent.putExtra("name", "張三")     // 車上捎帶"行李"（資料）
startActivity(intent)

// 方式二：隱式 Intent——只說要幹什麼，系統找誰幹
val web = Intent(Intent.ACTION_VIEW, Uri.parse("https://developer.android.com"))
startActivity(web)   // 系統會開啟瀏覽器/或讓你選
\end{lstlisting}

兩種方式各有用途：顯式用於自己 App 內部的頁面跳轉；隱式用於呼叫系統或其他 App 的能力（開網頁、打電話、分享）。Compose 專案裡，在 Composable 中用 LocalContext.current 拿到 Context 就能 startActivity。

\booksection{四、多頁面與導航}

真實 App 都有多個頁面。Compose 官方推薦的導航方案是 Navigation Compose（依賴庫 androidx.navigation）：

\begin{lstlisting}[language=Kotlin]
val navController = rememberNavController()

NavHost(navController, startDestination = "home") {
    composable("home") { HomeScreen(navController) }
    composable("detail/{name}") { entry ->
        val name = entry.arguments?.getString("name")
        DetailScreen(name ?: "未知")
    }
}

// 在頁面裡跳轉：
navController.navigate("detail/張三")     // 帶引數
navController.popBackStack()              // 返回上一頁
\end{lstlisting}

\begin{tcbinfo}{小貼士：頁面多了怎麼辦？}
專案長大後，把所有頁面路由集中寫在一個檔案裡（比如 AppNavGraph.kt），統一管理，別讓“跳轉到哪”的邏輯散落在各個頁面。這是從“能跑”走向“可維護”的第一步。

\end{tcbinfo}

\booksection{五、順帶認識：Fragment 是什麼（以及為什麼本書沒教）}

查老教學時你一定會撞見 Fragment——它是 Activity 裡的“可複用介面片段”，在老式多頁面架構裡是絕對主角。但 Compose 時代，頁面直接用 Composable 函式 + Navigation 路由就夠了，官方也已弱化 Fragment 的地位。本書因此沒把它放進主線。

給你的建議：新專案不主動用 Fragment；但讀老程式碼、面試被問到時，知道它是“Activity 裡嵌的介面片段”即可。學習資源永遠優先選 2023 年以後的——安卓技術棧的“老黃曆”更新得比想象中快。

\begin{tcbpractice}{來動手：給計數器裝上"記憶"}
把第 51 章的計數器改造一下：(1) 把 remember 換成 rememberSaveable，旋轉螢幕驗證數字還在；(2) 加一個“關於”按鈕，用 Navigation 跳到第二個頁面顯示版本號；(3) 在兩個頁面的 onCreate/onPause 裡列印日誌（Log.d），用 Logcat 觀察生命週期呼叫順序。做完這三步，生命週期就不再是紙上談兵。

\end{tcbpractice}

\begin{tcbsummary}{本章小結}
\begin{itemize}
  \item 四大元件：Activity（門面）、Service（後臺）、BroadcastReceiver（喇叭）、ContentProvider（資料櫃）
  \item 生命週期六回撥：onCreate onStart onResume onPause onStop onDestroy
  \item 旋轉螢幕會銷燬重建 Activity：用 rememberSaveable 保住狀態
  \item Intent 是"車票"：顯式跳自家頁面，隱式調系統能力
  \item 多頁面用 Navigation Compose：NavHost + composable 路由
\end{itemize}

\end{tcbsummary}
